% v14-cplus convergent evidence section: explicit mapping of L1-L6 to body
% material. Inserted between sec_mechanisms and sec_robustness in the
% paper_v14cplus.tex flow.

\section{Convergent Evidence}
\label{sec:convergent_evidence}

We assemble in this section the six lines of empirical evidence
introduced in the introduction (\S\ref{sec:intro}). Each line is
verifiable separately and corresponds to a material in this paper or
in its online appendix; we summarize the cross-references here so
that a reader can audit each line independently. None of the lines,
on its own, identifies a causal mechanism; jointly they constrain the
class of readings consistent with the data.

\paragraph{L1: Conditional within-item price premium.}
The OLS specifications of Equation~\eqref{eq:ols_main} yield a
conditional FL-presence price coefficient of $\valOLSGeneralPBU$ in
the within-item, within-year, within-PBU specification, with a
range $\valFLpremTopReg$--$\valFLpremBottomReg$ across modal
subsamples and specifications (Section~\ref{sec:results_ols},
Table~\ref{tab:prices}). The matching diagnostics in
Section~\ref{sec:robustness} (Table~\ref{tab:item_level_scope_match})
show the coefficient does not survive overlap-strict reweighting,
and we treat the gap as conditional, not causal.

\paragraph{L2: Oversight-capacity gradient.}
The conditional price coefficient varies by a factor of $\valPBUgradientRatio$
across PBU-size quartiles, monotonically: $\valOversightQone$ in Q1
(smallest PBUs), $\valOversightQtwo$ in Q2, $\valOversightQthree$ in
Q3, and $\valOversightQfour$ in Q4 (largest)
(Section~\ref{sec:robustness}, ``Oversight heterogeneity'').
Smaller PBUs are associated with several institutional features in
addition to monitoring capacity (item composition, market structure,
fiscal capacity); we do not interpret this gradient as a clean
causal estimate of monitoring effects, but the magnitude and
monotonicity are consistent with the framework's prediction
$\partial m^*/\partial \theta_k < 0$
(Proposition~\ref{prop:optimal_m}).

\paragraph{L3: Structural Regime~2 selection across modalities.}
Bid-distribution regime selection by BIC (Appendix~\ref{app:proofs},
Table~\ref{tab:structural_modal}) prefers Regime~2 over Regime~1
in both modalities by large margins: $\Delta\mathrm{BIC}=\valStructConvDeltaBIC$
in convite, $\valStructPregDeltaBIC$ in pregão. Pregão exhibits
sharper coordination ($\hat{\sigma}_c/\hat{\sigma}_g = \valStructPregRatio$
vs.\ $\valStructConvRatio$ in convite). Cover bidding in BEC is
therefore anchor-following (Regime~2), not detection-evading
through dispersion (Regime~1), in both modal regimes.

\paragraph{L4: Bayesian latent-class agreement.}
A two-component Gaussian mixture model on bid distributions
(Section~\ref{sec:classification},
Table~\ref{tab:lca_validation}) classifies $\valLCAcoverN$ firms as
the cover-class type. The mean posterior probability of cover-class
membership is $\valLCAprobCoverFL$ for FL firms versus
$\valLCAprobCoverNonFL$ for non-FL---a $\valLCAratio\!\!\times$ ratio.
Latent-class classification uses a different operationalization
from our construct (bid-distribution shape vs.\ participation-count
signal); the agreement is complementary, not strictly independent,
because both methods rely on the same underlying BEC bid-level
data.

\paragraph{L5: Repeated FL--winner pairs above null.}
Permutation of FL labels within the always-loser pool ($1{,}000$
iterations, uniform shuffle in
\texttt{scripts/45\_legacy\_m1m3\_perm\_welfare.R}) yields a null
mean of $\valDyadicPairsPerm$ FL--winner pairs sharing $\geq 5$
tender-items; the observed value is $\valDyadicPairsObs$, a
$\valDyadicRatio\!\!\times$ ratio ($p < 0.001$). The top-10 pairs
share $\valDyadicTopTenObs$ tender-items on average versus
$\valDyadicTopTenPerm$ under permutation
($\valDyadicTopTenRatio\!\!\times$,
Table~\ref{tab:dyadic_permutation}). Persistent FL--winner
adjacencies are informative even within the construct's own
population; the test does not distinguish coordination from
geographic or product-market co-location.

\paragraph{L6: CADE adjudication consistency check.}
The FL marker discriminates $\valCobidders$ adjudicated
CADE-cobidder firms within the always-loser stratum at AUC
$\valAUCFLfirm$ ($95\%$ CI $\valAUCFLfirmCI$;
Table~\ref{tab:horse_race_v14}). The construct does not
discriminate the $\valDirectCADE$ direct CADE defendants in the
broader BEC universe (AUC $\valAUCdirectStd$). The within-stratum
result is informative about the construct's behavior at the
cobidder boundary, not about cartel-detection scope at the firm
level (Section~\ref{sec:cade}).

\paragraph{Summary.}
Lines L1 and L2 measure the same conditional price-association
phenomenon at different cuts of the data and are not strictly
independent. Lines L3, L4, and L5 use bid-level information from
overlapping populations; their independence from L1 is partial.
Line L6 uses external enforcement labels and is structurally most
distinct, but its scope is limited to the always-loser stratum.
We therefore characterize the evidence as
\emph{multiple complementary lines} rather than fully independent
tests; the framework survives elimination of any one line, but the
lines reinforce each other rather than testing strictly orthogonal
implications.
